\section{Certificate ranks and movement in a fixed barrier metric}
\label{sec:exposed-movement}
This section bounds primal movement: the distance, in the Hessian
metric of a fixed barrier, from a reference point to the set of
$\epsilon$-accurate points. A lower bound on the barrier parameter
does not by itself bound this distance; one also needs to know which
part of the cone the objective forces toward its boundary. A certificate
for the objective describes this part, and the distance bound grows with
its total Jordan rank. The weighted exposed-minor estimate
(Theorem~\ref{thm:movement-minor}) is also proved in the companion
manuscript~\cite{CentralPathCompanion}; we include its proof to keep
this paper self-contained. The Jordan identities are
classical~\cite{FK1994}, as is the conversion from bounded local steps
to Riemannian distance~\cite{NT2002}. What is new is their combination
with the certificate-rank bounds of
Sections~\ref{sec:rank}--\ref{sec:orbits} for a fixed cone dictionary,
using the same certificate in both bounds.

For a barrier $F$ with positive-definite Hessian on its effective affine
domain $\Omega$, write
\[
 \norm{h}_{F,x}^2=D^2F(x)[h,h],\qquad
 d_F(x,y)=\inf_\gamma\int_0^1\norm{\dot\gamma(t)}_{F,\gamma(t)}\,dt.
\]
The infimum is over piecewise $C^1$ curves in $\Omega$ joining the two
points. For a set $T\subset\Omega$, put $d_F(x,T)=\inf_{y\in T}d_F(x,y)$,
with value $+\infty$ if $T$ is empty. Any free affine direction
along which the slack is constant is removed first.

\subsection{The weighted exposed-minor estimate}
Let $K=\prod_iK_i$ be a product of symmetric cones with their Jordan
trace inner products. Let $\Omega=\intr K\cap\mathcal A\ne\varnothing$
for an affine subspace $\mathcal A$. Suppose an affine objective $\ell$
has finite supremum $\ell_*$ and a genuine certificate $s=(s_i)\in K$
satisfying
\begin{equation}\label{eq:movement-certificate}
 g(x):=\ell_*-\ell(x)=\sum_i\ip{s_i}{x_i}\qquad(x\in\Omega).
\end{equation}
This identity must hold on the entire feasible affine slice. For each
nonzero $s_i$, let $c_i=\supp s_i$, $q_i=\jrank s_i$, and let
$\det_{c_i}$ denote the determinant in the Peirce algebra with unit $c_i$.
The expressions below omit factors with $q_i=0$.

\begin{theorem}[Weighted exposed-minor distance]\label{thm:movement-minor}
Use the restricted barrier
\[
 F(x)=-\sum_i\alpha_i\log\det_i x_i,\qquad \alpha_i\geq1,
 \qquad Q_\alpha=\sum_i\alpha_iq_i>0.
\]
For a fixed reference $x^0\in\Omega$, define
\begin{equation}\label{eq:movement-delta}
 \Delta_\alpha=Q_\alpha
 \left[\prod_{i:q_i>0}
 \frac{\{\det_{c_i}(s_i)\det_{c_i}(P(c_i)x_i^0)\}^{\alpha_i}}
      {\alpha_i^{\alpha_iq_i}}\right]^{1/Q_\alpha}.
\end{equation}
Then $\Delta_\alpha>0$ and, for every $\epsilon>0$,
\begin{equation}\label{eq:movement-minor-distance}
 d_F(x^0,\{x:g(x)\leq\epsilon\})
 \geq\sqrt{Q_\alpha}\,[\log(\Delta_\alpha/\epsilon)]_+.
\end{equation}
Here $[t]_+=\max(t,0)$. The bound applies to every feasible path in
$\Omega$, including paths that change coordinates not exposed by $s$
or additional lift variables such as completion variables. Neither
strict complementarity nor boundedness of those variables is assumed.
\end{theorem}
\begin{proof}
For an idempotent $c$ of rank $q$, put $a=P(c)x$ and
$\phi_c(x)=-\log\det_c a$. The element $a$ is interior to its Peirce
cone: $x\succeq\delta e$ implies $a\succeq\delta c$. The gradient and
inverse Hessian of $-\log\det x$ are $-x^{-1}$ and $P(x)$, respectively.
Since $P(c)$ is self-adjoint and $a^{-1}\in V(c,1)$,
\[
 D\phi_c(x)[h]=-\ip{a^{-1}}h,\qquad
 \norm{D\phi_c(x)}_{-\log\det,*}^2
 =\ip{a^{-1}}{P(x)a^{-1}}=q.
\]
The last equality follows from the fundamental identity
\[
 P(c)P(x)P(c)=P(P(c)x)=P(a),\qquad
 P(a)a^{-1}=a,\qquad \ip{a^{-1}}a=q.
\]
These are Jordan identities, so they also hold in the exceptional algebra.
For $\Phi(x)=\sum_{i:q_i>0}\alpha_i\phi_{c_i}(x_i)$, the full product
metric consequently gives $\norm{D\Phi}_{F,*}^2=Q_\alpha$.
Restricting a covector and the metric to an affine tangent subspace can
only decrease its dual norm. Hence
\begin{equation}\label{eq:movement-potential-lipschitz}
 |D\Phi(x)[h]|\leq\sqrt{Q_\alpha}\norm{h}_{F,x}
 \quad(h\in T_x\Omega).
\end{equation}

For a feasible endpoint $x$, set $g_i=\ip{s_i}{x_i}>0$ on the active
factors. In $V(c_i,1)$, the positive element
$P(s_i^{1/2})P(c_i)x_i$ has trace $g_i$ and determinant
$\det_{c_i}(s_i)\det_{c_i}(P(c_i)x_i)$. Spectral AM--GM therefore gives
\[
 \det_{c_i}(P(c_i)x_i)
 \leq\det_{c_i}(s_i)^{-1}(g_i/q_i)^{q_i}.
\]
If $\sum_i g_i\leq\epsilon$, maximizing
$\sum_i\alpha_iq_i\log(g_i/q_i)$ over these positive $g_i$ gives
$g_i=\epsilon\alpha_iq_i/Q_\alpha$. It follows that
\[
 \prod_i\det_{c_i}(P(c_i)x_i)^{\alpha_i}
 \leq (\epsilon/Q_\alpha)^{Q_\alpha}
 \prod_i\alpha_i^{\alpha_iq_i}\det_{c_i}(s_i)^{-\alpha_i}.
\]
Comparison with the reference point yields
$\Phi(x)-\Phi(x^0)\geq Q_\alpha\log(\Delta_\alpha/\epsilon)$.
Integrating~\eqref{eq:movement-potential-lipschitz} along any feasible
curve proves the result.
\end{proof}

\begin{corollary}[Bounded-move consequence]\label{cor:movement-rounds}
Suppose a procedure starts at $x^{\rm in}$ with
$d_F(x^{\rm in},x^0)\leq D_0$ and makes $T$ rounds, each with metric
displacement at most $B_0>0$, before attaining gap at most $\epsilon$.
Then
\begin{equation}\label{eq:movement-round-count}
 T\geq
 \frac{[\sqrt{Q_\alpha}\log(\Delta_\alpha/\epsilon)-D_0]_+}{B_0}.
\end{equation}
In particular, for an integer $m\geq1$ and $0<R<1$, if each round
consists of at most $m$ feasible forward chords $x+th$,
$0\leq t\leq1$, with starting local norm $\norm h_{F,x}\leq R$, one may
take $B_0=-m\log(1-R)$. If $\norm{x^{\rm in}-x^0}_{F,x^0}\leq\rho<1$,
one may take $D_0=-\log(1-\rho)$. More generally, if the procedure
makes $M\geq1$ chords with starting local norms $r_j<1$ and
$D=[\sqrt{Q_\alpha}\log(\Delta_\alpha/\epsilon)-D_0]_+$, then
\begin{equation}\label{eq:movement-variable-chords}
 \sum_{j=1}^M-\log(1-r_j)\geq D,\qquad
 \max_j r_j\geq1-e^{-D/M}.
\end{equation}
\end{corollary}
\begin{proof}
The triangle inequality proves~\eqref{eq:movement-round-count}.
For a chord $x+th$, the standard self-concordant Hessian comparison gives
$\norm h_{F,x+th}\leq\norm h_{F,x}/(1-t\norm h_{F,x})$.
Its integral is at most $-\log(1-\norm h_{F,x})$. Summing these
bounds proves the variable-chord assertion as well; see also
\cite[Lemma~3.2 and Corollary~3.2]{NT2002}.
\end{proof}

The square-root coefficient is sharp on the full product cone. A
quadratic-representation isometry sends $x^0$ to the unit. In a Jordan
frame diagonalizing the transformed $s_i$, multiply exactly its $q_i$
positive-support coordinates by $t$ and keep all remaining coordinates
equal to one. This curve has gap $t\sum_i\ip{s_i}{e_i}$ and length
$\sqrt{Q_\alpha}\log(1/t)$ from the unit. Thus, on the full cone,
\begin{equation}\label{eq:movement-full-cone-sharp}
 \lim_{\epsilon\downarrow0}
 \frac{d_F(x^0,\{g\leq\epsilon\})}{\log(1/\epsilon)}
 =\sqrt{Q_\alpha}.
\end{equation}
In this statement the certificate is transformed with the coordinates;
its pairing, and hence the gap, is unchanged. An affine slice can exclude
the displayed curve, so~\eqref{eq:movement-full-cone-sharp} does not
assert a matching upper bound for every slice.

Hauser and G\"uler's classification
\cite[Theorem~5.5]{HauserGuler2002} says that every self-scaled barrier
on a symmetric cone is, up to an additive constant, of the displayed
form on its irreducible factors, with $\alpha_i\geq1$. Consequently the
theorem covers all self-scaled barriers on a fixed symmetric cone.
Changing the weights changes both $Q_\alpha$ and $\Delta_\alpha$; the
bound must use $\Delta_\alpha$, not its unit-weight value.
The classification gives no analogous assertion for arbitrary coupled
barriers on an affine slice.

\subsection{Movement bounds from certificate ranks}
For unit weights, put $Q=\sum_iq_i$ and write $\Delta$ for
\eqref{eq:movement-delta}. For a fixed lift, reference, objective, and
certificate, Theorem~\ref{thm:movement-minor} gives
\begin{equation}\label{eq:movement-rank-limit}
 \liminf_{\epsilon\downarrow0}
 \frac{d_F(x^0,\{g\leq\epsilon\})}{\log(1/\epsilon)}
 \geq\sqrt Q.
\end{equation}
The following corollary applies this bound to the earlier rank results.
\begin{corollary}[Dictionary consequences]\label{cor:movement-dictionary}
Let the fixed symmetric-cone dictionary have capacity $B$ from
\eqref{eq:dictionary-capacity}, and use its standard product barrier on
the relative-Slater lift.
For each lift of $\B_2^s$, on a dense open full-measure set of supports,
the right-hand side of~\eqref{eq:movement-rank-limit} is at least
$\sqrt{\lceil(s-1)/B\rceil}$.
For every lift of $\prod_a\B_2^{s_a}$, there exist directions $u_a$
and genuine pure-row certificates $Y^a$ such that every fixed choice
of positive weights $\lambda_a$ gives an objective
$\ell(x)=\sum_a\lambda_a\ip{u_a}{\pi_a x}$ for which the
right-hand side is at least
\[
 \sqrt{\sum_a\left\lceil\frac{s_a-1}{B}\right\rceil}.
\]
Under the globally bi-$C^1$ genuine-selection hypotheses of
Theorem~\ref{thm:smooth-rank}, the right-hand side is instead at least
$\sqrt{\kappa_{\mathscr K}(s)}$ for some support of a single ball.
\end{corollary}
\begin{proof}
For one ball choose any genuine certificate at a support in
Theorem~\ref{thm:rank-frontier}'s regular set. For the product take
exactly the certificate $s_i=\sum_a\lambda_aY_i^a$ supplied by
Theorem~\ref{thm:product-existence}. It represents the full gap
$\sum_a\lambda_a(1-\ip{u_a}{\pi_a x})$ and has the asserted aggregate
rank. In the smooth case use the genuine selected certificate whose
rank is bounded below in Theorem~\ref{thm:smooth-rank}. Apply
\eqref{eq:movement-rank-limit} in each case.
\end{proof}
The product assertion uses only the existence of one such aggregate
certificate; it does not say that every aggregate certificate has this
rank, and it does not minimize over the whole dual fiber. One high-rank certificate
suffices because its determinant estimate holds at every accurate primal
endpoint. For self-scaled barriers the same statements hold with the
weighted $Q_\alpha\geq Q$. Since $\Delta_\alpha$ can depend on the lift,
objective, reference, and certificate, the corollary is an asymptotic
statement for a fixed instance, not a finite-accuracy estimate uniform
over all lifts. In particular, this argument does not allow replacing
$Q$ by a larger ambient or restricted barrier parameter.

\subsection{Affine PSD pencils and spectral scales}
A second form of the estimate compresses an affine matrix pencil $L$
to a subspace. It is useful when an exposed minor $\det(R^*LR)$ is
known directly, including in completion formulations. The convexity
used is classical matrix-fractional convexity
\cite[Example~3.4]{BoydVandenberghe2004}.
The matrix-ball specialization and its rank-profile interpretation
also appear in~\cite{CentralPathCompanion}.

\begin{proposition}[Principal-minor contraction]\label{prop:movement-psd-compression}
Let $L(x)=L_0+\mathcal A x\succ0$ be a real symmetric or complex
Hermitian affine pencil and $F=-\log\det L$ on an effective affine
domain with positive-definite Hessian. For an injective matrix $R$ with
$r\geq1$ columns put $\psi_R(x)=-\log\det(R^*L(x)R)$. Then
\begin{equation}\label{eq:movement-psd-contraction}
 D^2\psi_R\preceq D^2F,\qquad
 |D\psi_R(x)[h]|\leq\sqrt r\norm h_{F,x}.
\end{equation}
If the compressed function has the sharper restricted gradient bound
$|D\psi_R[h]|\leq\sqrt\vartheta\,(D^2\psi_R[h,h])^{1/2}$ for
$\vartheta>0$, the
second constant can be replaced by $\sqrt\vartheta$. Consequently
\begin{equation}\label{eq:movement-psd-ratio}
 d_F(x^0,x)\geq\frac1{\sqrt\vartheta}
 \left[\log\frac{\det(R^*L(x^0)R)}{\det(R^*L(x)R)}\right]_+,
\end{equation}
with $\vartheta=r$ always permitted.
\end{proposition}
\begin{proof}
First take $R$ isometric and complete it to a unitary basis. In that
basis write $L=\left(\begin{smallmatrix}A&B\\B^*&C\end{smallmatrix}\right)$,
where $A=R^*LR$. If $R$ is square there is nothing to prove. Otherwise
\[
 F-\psi_R=-\log\det(C-B^*A^{-1}B).
\]
The Schur complement is matrix concave because the matrix-fractional
map $B^*A^{-1}B$ is matrix convex: test the scalar convex function
$(Bv)^*A^{-1}(Bv)$ against each fixed vector $v$. The scalar
matrix-fractional convexity follows directly from its Schur-complement
epigraph, over either field. Composition with the convex,
Loewner-decreasing function $-\log\det$ is convex. This proves the
Hessian inequality. Also, with $T=A^{-1/2}(D A[h])A^{-1/2}$,
\[
 D\psi_R[h]=-\tr T,\qquad D^2\psi_R[h,h]=\tr(T^2),\qquad
 |\tr T|\leq\sqrt r\,\sqrt{\tr(T^2)}.
\]
A thin QR factorization of an arbitrary injective $R$ changes
$\psi_R$ only by a constant, proving the general case. Integration
proves~\eqref{eq:movement-psd-ratio}, also with any stated sharper
constant.
\end{proof}
For example, for the matrix-ball pencil
$L(X)=\left(\begin{smallmatrix}I_p&X\\X^*&I_q\end{smallmatrix}\right)$,
compression to $1\leq r\leq\min(p,q)$ chosen left coordinates and all right
coordinates has determinant $\det(I_r-U^*XX^*U)$, where $U$ is the
$p$ by $r$ matrix of the selected orthonormal left coordinate vectors.
Its restricted gradient parameter
is $r$, although the compressed pencil has order $r+q$, so one should
use $\vartheta=r$. The gradient bound remains
valid when this compressed Hessian is singular in some directions.

\begin{theorem}[Exposed spectral profile]\label{thm:movement-profile}
In Proposition~\ref{prop:movement-psd-compression}, let $W\succeq0$
have rank $r\geq1$, and suppose the target satisfies
$\tr(WL(x))\leq\epsilon$. Write $W=QDQ^*$ with $Q^*Q=I_r$ and
$D\succ0$, and set
\[
 Z_0=D^{1/2}Q^*L(x^0)QD^{1/2},\qquad
 d_1\geq\cdots\geq d_r>0
\]
for its eigenvalues. Then
\begin{equation}\label{eq:movement-profile}
 d_F(x^0,x)\geq
 \max_{1\leq m\leq r}\sqrt m
 \left[\log\frac{m(d_1\cdots d_m)^{1/m}}{\epsilon}\right]_+.
\end{equation}
The same bound divided by $-\log(1-R_0)$ is a lower bound on the
number of feasible forward chords of local norm at most $0<R_0<1$.
\end{theorem}
\begin{proof}
Take top-$m$ orthonormal eigenvectors $U_m$ of $Z_0$ and use the
injective compression $R_m=QD^{1/2}U_m$. Its reference determinant is
$\prod_{j\leq m}d_j$. At the endpoint,
\[
 \tr(R_m^*L(x)R_m)\leq\tr(WL(x))\leq\epsilon,
 \qquad
 \det(R_m^*L(x)R_m)\leq(\epsilon/m)^m.
\]
Apply Proposition~\ref{prop:movement-psd-compression} with
$\vartheta=m$, then maximize over $m$. The chord statement follows
from Corollary~\ref{cor:movement-rounds}'s chord estimate.
\end{proof}
For $m=r$ this is the unweighted exposed-minor estimate. Taking $m<r$
can avoid a large loss when some positive eigenvalues are tiny.
The quantities $d_i$ include both the exposing matrix and the reference
slack; if $L(x^0)=I$, they are precisely the positive eigenvalues of $W$.

Denote the right side of~\eqref{eq:movement-profile} by
$\mathcal E_r(\epsilon)$ and put $\Lambda=\log(1/\epsilon)$. For
two eigenvalue profiles we give asymptotics of this lower bound; we do not
claim they equal the exact distance for an arbitrary pencil.
For $d_i=e^{-a(i-1)}$, $a>0$, and any fixed $\delta>0$, if
$r\geq(2/(3a)+\delta)\Lambda$ for all sufficiently small
$\epsilon$, then
\begin{equation}\label{eq:movement-geometric-profile}
 \mathcal E_r(\epsilon)
 =\left(\frac23\sqrt{\frac{2}{3a}}+o(1)\right)\Lambda^{3/2}.
\end{equation}
For $d_i=i^{-\beta}$, $\beta>1$, put
$C_\beta=e^{(2-\beta)/(\beta-1)}$. If
$r\geq(C_\beta+\delta)\epsilon^{-1/(\beta-1)}$, then
\begin{equation}\label{eq:movement-polynomial-profile}
 \mathcal E_r(\epsilon)
 =\left(2(\beta-1)\sqrt{C_\beta}+o(1)\right)
       \epsilon^{-1/[2(\beta-1)]}.
\end{equation}
To verify~\eqref{eq:movement-geometric-profile}, maximize
$\sqrt m[\Lambda+\log m-a(m-1)/2]_+$ over positive real $m$.
Its positive maximizer satisfies
\[
 m_*=\frac{2}{3a}(\Lambda+\log m_*+a/2+2)
     =\frac{2\Lambda}{3a}+O_a(\log\Lambda).
\]
The bracket there is $am_*-2$. Rounding $m_*$ to an integer gives
\eqref{eq:movement-geometric-profile}. For the polynomial case,
Stirling's estimate gives
\[
 \log\frac{m(m!)^{-\beta/m}}{\epsilon}
 =\Lambda-(\beta-1)\log m+\beta+O_\beta(\log(m+1)/m).
\]
The error after multiplication by $\sqrt m$ is bounded uniformly in
$m\geq1$. The corresponding continuous expression is maximized at
$m_*=C_\beta\epsilon^{-1/(\beta-1)}$, with bracket
$2(\beta-1)$, proving~\eqref{eq:movement-polynomial-profile}.
In both cases integer rounding is a lower-order error. If the rank $r$ is
smaller than this scale, the maximum over $m\leq r$ in
\eqref{eq:movement-profile} is the applicable bound. In particular,
for fixed finite $r$, its leading behavior as $\epsilon\downarrow0$
is $\sqrt r\log(1/\epsilon)$.

The orders of these lower bounds are attained by particular diagonal
pencils. For each finite list $d_i$, take
\[
 L(q)=\operatorname{Diag}(q_1,\ldots,q_r,2-q_1,\ldots,2-q_r),
 \quad q\in(0,2)^r,\quad q^0=\mathbf1,
 \quad W=\operatorname{Diag}(d_1,\ldots,d_r,0,\ldots,0).
\]
The target is $\sum_i d_iq_i\leq\epsilon$. On $0<q\leq1$ define
\[
 \varrho(q)=\int_q^1\sqrt{t^{-2}+(2-t)^{-2}}\,dt,
 \qquad
 \log(1/q)\leq\varrho(q)\leq\sqrt2\log(1/q).
\]
The coordinates $\varrho(q_i)$ turn the product metric into the Euclidean
metric. Set $q_i^*=\min\{1,\epsilon/(rd_i)\}$. This endpoint has
gap at most $\epsilon$, and a straight segment in the $\varrho$
coordinates gives a feasible path of length at most
\begin{equation}\label{eq:movement-diagonal-upper}
 \sqrt2\left(\sum_{i=1}^r[\log(rd_i/\epsilon)]_+^2\right)^{1/2}.
\end{equation}
Choose $r=\lceil c\Lambda\rceil$ with fixed $c>2/(3a)$ in the
geometric case. Every summand in~\eqref{eq:movement-diagonal-upper}
is $O_a(\Lambda^2)$, so the upper order is $\Lambda^{3/2}$.
In the polynomial case take
$r=\lceil c\epsilon^{-1/(\beta-1)}\rceil$, $c>C_\beta$.
Then $\log(rd_i/\epsilon)$ is a bounded constant minus
$\beta\log(i/r)$, whose positive-part square has a finite integral
on $(0,1]$. The sum is $O_\beta(r)$, giving the matching order
$\epsilon^{-1/[2(\beta-1)]}$. Constants here may depend on the
fixed $c$ and decay parameter. These examples use arbitrary
feasible paths, not central paths, in specially chosen box spectrahedra
whose rank grows at the specified rate; they give no matching upper
bound for general affine pencils.
